\section{Del entrenamiento inicial al refinamiento del método}
\label{sec:observaciones-entrenamiento}
El SFT inicial muestra una separación entre la pérdida de entrenamiento y la de validación en las últimas épocas. Este comportamiento justifica observar ambos conjuntos al revisar la duración del ajuste, en lugar de seguir únicamente la pérdida calculada durante las actualizaciones.

En el piloto posterior, la caída de recompensa al retirar las ayudas señala otra dificultad: obtener una respuesta bien puntuada cuando se proporcionan conceptos no equivale a recuperarlos a partir de la consulta. La validación sin ayudas permite seguir esa diferencia durante todo el recorrido, incluso en las etapas donde la recompensa de entrenamiento permanece alta.

La evaluación inicial de respuestas también motivó revisar qué se exige en cada rúbrica. La segunda etapa vincula los criterios con las obligaciones de la pregunta y conserva evidencia de la referencia, buscando evitar que detalles secundarios se conviertan en requisitos. El capítulo~\ref{ch:resultados-adaptacion} examina las respuestas mediante HealthBench y distingue ese análisis del contraste con el trabajo externo.
